\section{Related Work}

Gridded population mapping has developed from simple areal representations toward increasingly sophisticated methods for redistributing aggregated population counts into regular grid cells. A large body of work now distinguishes between products that primarily preserve counts within source administrative units, products that dasymetrically redistribute population using ancillary spatial data, and model-based products that combine census counts with geospatial covariates to infer finer-scale spatial variation \cite{Leyk2019,Mennis2009}. Across these product families, the core methodological challenge remains the same: population is usually observed only in aggregated reporting units, whereas many analytical applications require spatially explicit estimates on user-defined grids.

A major branch of this literature focuses on top-down disaggregation, in which official or census-based totals are redistributed across grid cells using weighting layers derived from land cover, settlement extent, buildings, infrastructure, or other ancillary variables. In this sense, the problem is closely related to dasymetric mapping, where ancillary spatial evidence is used to constrain where population can plausibly be allocated within a reporting unit \cite{Mennis2009,Stevens2015}. These approaches have become especially important for disaster exposure analysis, public-health planning, and urban comparison because they provide flexible population surfaces that can be aggregated to arbitrary spatial units while remaining linked to official totals at coarser scales \cite{Leyk2019,WorldPopMethods}.

Among the most widely used model-based products, WorldPop has established a prominent top-down framework that combines administrative population counts with geospatial covariates and machine-learning-based weighting surfaces, including random-forest-based redistribution workflows in multiple regional and global products \cite{Stevens2015,Sorichetta2015,Lloyd2017}. WorldPop itself emphasizes that gridded population estimates are highly useful for flexible spatial analysis but do not replace the need for a full census, particularly where richer demographic and socioeconomic detail is required \cite{WorldPopMethods}. This distinction is important for the present study because it reinforces the difference between a calibrated gridded proxy surface and true cell-level census observation.

A second major line of work is represented by the Global Human Settlement Layer population products. GHS-POP disaggregates harmonized census or administrative population counts to grid cells using built-up information from the GHSL framework, and the associated GHS-POP2G tools operationalize this logic across multiple spatial resolutions \cite{Freire2016,GHSL2023,GHSPOP2G2020}. Compared with more heavily model-driven approaches, this family of methods places strong emphasis on transparent and scalable redistribution workflows tied to built-up structure. For the purposes of this paper, GHS-POP is therefore relevant both as an external benchmark and as an example of officially anchored gridded population allocation rather than cell-level truth recovery.

Recent work has also explored higher-resolution and more localized population mapping strategies using combinations of building data, roads, points of interest, remote sensing, and other geospatial signals. Some studies use machine-learning approaches at local or regional scales, while others emphasize open-data or building-based disaggregation workflows designed for settings where conventional fine-resolution labels are unavailable \cite{Qiu2020,Pirowski2024,Popcorn2024}. Although these approaches differ in inputs and modeling assumptions, they share an important practical insight: in data-scarce environments, useful intra-urban population surfaces are often produced not by direct observation of true cell-level counts, but by transparent proxy modeling constrained by coarser official totals.

An equally important theme in the literature concerns validation and fitness for use. Large-scale gridded population products differ substantially in source inputs, assumptions, resolution, temporal alignment, and intended use, which means that no single product or validation statistic is universally sufficient across applications \cite{Leyk2019}. Empirical evaluations have shown that local performance can vary materially by setting and that deprived or complex urban environments remain particularly challenging for gridded population estimation \cite{Thomson2021,Yin2021}. This creates a methodological gap for studies in data-scarce urban contexts: researchers often need a reproducible baseline whose limitations are explicit and whose evaluation is multi-layered rather than dependent on a single benchmark.

City1 v2 is positioned within this literature as a \emph{baseline system contribution}. It does not claim best-possible truth recovery and does not present the resulting grid as true census population at the cell level. Instead, it contributes a reproducible open-data pipeline for intra-urban proxy population surface generation, calibrated by official city totals and evaluated through a multi-level validation stack tailored to data-scarce urban settings. In this sense, the contribution is closer to a disciplined, transparent, and bounded baseline for intra-urban analysis than to a claim of definitive population reconstruction.
